

%%%%%%%%% FIGURE MARGINS: %%%%%%%%%%%%%  
\makeatletter
\renewenvironment{figure}[1][]{%
  \def\@captype{figure}%
  \par\centering}
  {\par}
\makeatother

%%%%%%%%%%%%%%%%%%%%%%%%%%%%%%
%% Caption width == figure width
% https://tex.stackexchange.com/questions/53497/how-do-i-control-the-exact-dimension-of-a-picture-in-latex
% https://tex.stackexchange.com/questions/53497/how-do-i-control-the-exact-dimension-of-a-picture-in-latex

\usepackage{printlen}
% \usepackage{caption}
\usepackage[font=small]{caption}
%  %% CAPTION:
%     \node (capt) [yshift=11mm,below=of A,text width=\imagewidth, font=\small, align=justify]%\scriptsize , align=center
%     {\figurecaptionfont #3};%\subframetitlefont
\newlength\imageheight
\newlength\imagewidth


%% maximum image height, which could be used to ensure same dims.
% Macros for reconizing the aspect ratio:
\makeatletter
\ifthenelse{%
  \lengthtest{\beamer@paperwidth=16cm}}% condition (16cm is the width of the frame with an aspect ratio = 16:9)
    {% 16:9 aspect ratio
    \newcommand\maxfigheight{55mm} 	% 118 mm
    \newcommand\maxfigwidth{140mm} 	% 118 mm
    }
    {% 4:3 aspect ratio
    \newcommand\maxfigheight{55mm} 	% 118 mm
    \newcommand\maxfigwidth{100mm} 	% 118 mm
    }%do nothing
\makeatother
% 4:3 aspect ratio
% \renewcommand\framenumXpos{98mm}



%% Make 1 centered figure: 
\newlength{\mylength}   % for raising the figs
\newlength{\defaultYspace}   % for raising the figs 
\newcommand{\figOne}[3]
{\vspace{\defaultYspace}%default vspace
\begin{figure} 
  \settowidth\imagewidth{\includegraphics[#1]{#2}}
  \includegraphics[#1]{#2}
  \captionsetup{width=\the\imagewidth} %\the\imageheight
  \caption{#3}
 \end{figure}
 \settototalheight{\mylength}{#3}% calculate caption height
 %  % Raise figure if enough cites are mentioned:
 \ifPositif{\yshiftDueCites}
    {\raisebox{\yshiftDueCites+\mylength+3em}{}% Raise lists
    }{\raisebox{\mylength+2em}{}% Raise figure/caption up
     }
}

 

%%% Newest figure inserting command %%%
\newlength\Maximagewidth
\newlength\Maximageheight

\usepackage{pgfkeys}
%%%  key-value-based figures:
% \pgfkeys{
%     /fig/.is family, /fig,
%     default/.style =
%     {height = \maxfigheight,
%     width = \maxfigwidth,
%     capwidth = 0cm},
%     height/.estore in = \fig@height,
%     width/.estore in = \fig@width,
%     capwidth/.estore in = \fig@capwidth,
%     }
    
% noexpand magic?, the input arguments won't get exapnded?
% \newlength\test
% \newlength\dappa
\pgfkeys{
    /figu/.is family, /figu,
    default/.style =
    {height = 6cm,
    width = 8cm,
    m = 0cm,
    broadencaption = 0cm,},
    height/.store in = \figu@height,
    width/.store in = \figu@width,
    m/.store in = \figu@m,
    broadencaption/.store in = \figu@broadencaption,
    }

%% make e.g. width and height comparisons to be compatible with given input heights etc. 
%% would be also nice to be fully compatible with the \includegraphics[]{} functionality.
% noexpand magic?, the input arguments won't get exapnded?
%% difference between store and estore...?


%% NEW TIKZ BASED FIGURE PLACING COMMAND:
% a lot clearer to make separate macros to
% a) calculate the optimal figure size
% b) position the figures etc.

\newlength{\myheight}       % for raising the figs
\newlength{\captionheight}   % for raising the figs
\newlength{\tmpwidth}   % for raising the figs
\newlength{\figureYoffset}
\newcommand*{\yshift}{}
\pgfkeys{
    /figut/.is family, /figut,
    default/.style =
    {height = 6cm,
    width = 8cm,
    yshift = 0cm,
    broadencaption = 0cm, },
    height/.estore in = \figut@height, %
    width/.estore in = \figut@width, %
    yshift/.estore in = \yshift, %
    broadencaption/.estore in = \figut@broadencaption, %
    }

% \newcommand\maxfigheight{55mm} 	% 118 mm
\newlength{\maxfigheightnu}
\setlength{\maxfigheightnu}{60mm} %55mm
\newcommand{\fignu}[3][]{%
    \pgfkeys{/figut, default, #1}%
    %{
    % OFFSET USING CAPTION HEIGHT:
    % \settototalheight{\captionheight}{
    %     \begin{tikzpicture} 
    %     \node (capt) [yshift=11mm, text width=\imagewidth, font=\small, align=justify]
    %         {\subframetitlefont #3};
    %     \end{tikzpicture}
    % }
    \setlength{\tmpwidth}{\maxfigwidth}%
    \addtolength{\tmpwidth}{\figut@broadencaption}%
    \settoheight{\captionheight}{
        \begin{tikzpicture} 
        \node (capt) [yshift=11mm, text width=\tmpwidth, font=\small, align=justify]%\imagewidth
            {\subframetitlefont #3};
        \end{tikzpicture}
    }
    % \the\maxfigheightnu \\
    % \the\captionheight \\
    % \the\maxfigheightnu \\
    \addtolength{\maxfigheightnu}{-\captionheight}%\maxfigheightnu
    % \the\maxfigheightnu
    % OPTIMIZE FIGURE WIDTH / HEIGHT:
    \settowidth\Maximagewidth{\includegraphics[width=\maxfigwidth]{#2}}%
    \settowidth\Maximageheight{\includegraphics[height=\maxfigheightnu]{#2}}%
    \ifthenelse{\lengthtest{\Maximagewidth > \Maximageheight}}
        {% extract figure width for caption:
        \settowidth\imagewidth{\includegraphics[height=\maxfigheightnu]{#2}}
         }{
        % extract figure width for caption:
        \settowidth\imagewidth{\includegraphics[width=\maxfigwidth]{#2}}
        }
    
    % broaden \imagewidth to widen captions:
    \addtolength{\imagewidth}{\figut@broadencaption}%
    %
    %% PICTURE:
    % offset by half of the caption height
    \addtolength{\figureYoffset}{0.5\captionheight}
    % add slight downward offset for balanced slides [-0.2cm]
    % \addtolength{\figureYoffset}{-0.1cm}
    \addtolength{\figureYoffset}{-0.4cm} % for kolumbus fonts
    % optional handle to shift the figure & caption
    \addtolength{\figureYoffset}{\yshift}
    %
    \begin{tikzpicture} %%%% TIKZ %%%%
        [remember picture,overlay]
        % offset with \figureYoffset because of captions
        \node (A) [xshift=0mm,yshift=\figureYoffset, anchor=center] at (current page.center)
        {%% re-size the figure depending on above lengths:
          \ifthenelse{\lengthtest{\Maximagewidth > \Maximageheight}}{
                \includegraphics[height=\maxfigheightnu]{#2}
              }{
                % make figure
                \includegraphics[width=\maxfigwidth]{#2}
              }
          };
    %% CAPTION:
    \node (capt) [yshift=11mm,below=of A,text width=\imagewidth, font=\small, align=justify]%\scriptsize , align=center
    {\figurecaptionfont #3};%\subframetitlefont
    \end{tikzpicture} %%%% TIKZ %%%%
    %}
}


 
%% Make 2 figures side by side with separate captions, 1st to left, 2nd to right: 
\usepackage{etoolbox}% http://ctan.org/pkg/etoolbox
% \def\sA{1}
% \def\sB{2}
% \ifnumcomp{\sA}{<}{\sB}{ture}{false}

%%%%%%%%%%%%%%%%%%%%%%%%%%%%%%
%% INSERT FIGURE TO CUSTOM POSITION:
\newcommand{\figToArbPos}[6]
{\begin{tikzpicture}[remember picture,overlay]
    \node[xshift=#1,yshift=#2, anchor=#6] at (current #3) {\includegraphics[#4]{#5}};
\end{tikzpicture}}

%% Figures with caption to arbitrary position
\usetikzlibrary{positioning}
\newcommand{\figArbPosCap}[7]{
\begin{tikzpicture}
[remember picture,overlay]
    \node (A) [xshift=#1,yshift=#2, anchor=#6] at (current #3) {\includegraphics[#4]{#5}};
    \node (capt) [yshift=11mm,below=of A,text width=\linewidth, align=center,font=\scriptsize]%\scriptsize,font=\scriptsize
    { #7}; %\subframetitlefont
\end{tikzpicture}}


%%%%%%%%%%%%%%%%%%%%%%%%%%%%%%
%% TIKZ TO CUSTOM POSITION:
\newcommand{\tikzToArbPos}[1]
{\begin{tikzpicture}[remember picture,overlay]
    %\node[xshift=#1,yshift=#2] at (current #3) 
    #1 %{}
 \end{tikzpicture}
}

%% let's see if this works here.
\newcommand{\figArbPosCapB}[5]{
\begin{tikzpicture}
[remember picture,overlay]
    \node (A) [xshift=#1,yshift=#2, anchor=south] at (current page.center) {\includegraphics[#3]{#4}};
    \node (capt) [yshift=12mm,below=of A,text width=\linewidth, align=center,font=\small]%\scriptsize
    {\subframetitlefont #5};
\end{tikzpicture}}
